\clearpage\chapter{Random matrix theory}

\section*{1. Overview}
\subsection*{Intuitive}
\paragraph{The problem}
Suppose a data set has thousands of measurements and thousands of variables.
Its covariance matrix is then a huge table.  Even if the measurements contain
no real pattern, the table will have eigenvalues, large singular values, and
apparently important directions.  Which features are genuine, and which are
created by randomness?

The same question appears in physics.  A complicated quantum system may have
too many interacting particles for its exact energy levels to be calculated.
Instead of predicting every level, we can ask about the typical spacing of
nearby levels.  Random matrix theory was developed to answer such questions
about large matrices without pretending that every entry is known.


\subsection*{Formal}
An $N\times N$ random matrix is a matrix whose entries, or whose eigenvalue distribution, are governed by a specified probability law. Important observables include the eigenvalues $\lambda_1,\ldots,\lambda_N$, the empirical spectral measure $N^{-1}\sum_i\delta_{\lambda_i}$, and the normalized trace $N^{-1}\operatorname{Tr}f(M)$. The theory studies their limiting behavior as $N\to\infty$, with the ensemble and normalization stated explicitly.
\section*{2. History}
\subsection*{Historical development}
\begin{historylist}
\paragraph{History and the central idea}
Eugene Wigner introduced random matrices into nuclear physics in the 1950s.
He argued that the unknown interactions inside a heavy nucleus could be
modelled by a symmetric random matrix.  The individual entries were not
supposed to be literally random forces; randomness represented incomplete
knowledge and complicated interactions.  Freeman Dyson later organised the
main symmetry classes, and many mathematicians developed precise limit
theorems for eigenvalues and singular values \cite{random_matrix_theory_ref1}.

Let $X$ be an $m\times n$ matrix whose entries are independent measurements,
and form the sample covariance matrix
$S=(1/n)XX^{\mathsf T}$.  The eigenvectors of $S$ are the directions found by
principal component analysis.  If the entries of $X$ are only noise, the
eigenvalues are nevertheless not all equal.  For large $m$ and $n$, their
histogram approaches a predictable shape, the Marchenko--Pastur distribution,
when the ratio $m/n$ is held fixed.  A very large eigenvalue outside this
noise bulk is evidence for a signal, not merely evidence that a large matrix
always has a largest eigenvalue \cite{random_matrix_theory_ref2}.

\end{historylist}
\section*{3. Theory}
\subsection*{Core theory}
\paragraph{What the theory measures}
\bookfigure{random_matrix_spectrum}{Random matrix theory separates a typical noise spectrum from a possible signal outlier.}
There are three different questions, and confusing them leads to bad
interpretations.  The global eigenvalue distribution asks where most
eigenvalues lie.  The extreme-eigenvalue question asks how large the largest
one is.  The local question asks how close neighbouring eigenvalues are.

The answer depends on scaling and symmetry.  For a symmetric matrix with
entries of variance one, eigenvalues typically have size about the square
root of the matrix dimension.  Dividing by that scale produces a stable
limiting picture.  For a covariance matrix, the appropriate scale is set by
the number of observations and variables.  A theorem about one scaling
cannot be transferred to another without checking the assumptions.

\section*{4. Important theorems and results}
\begin{enumerate}[leftmargin=*,itemsep=0.35em]
\item \textbf{Marchenko--Pastur law.} For a large sample covariance matrix with
independent, suitably scaled entries, the empirical distribution of
eigenvalues approaches a deterministic limiting distribution.  In practice,
this gives a noise bulk against which unusually large principal components
can be compared.

\item \textbf{Wigner semicircle law.} For a large real symmetric random matrix with
centred entries of the appropriate variance, the rescaled eigenvalue
histogram approaches a semicircle.  The theorem explains why the detailed
distribution of individual entries often matters less than the matrix's
symmetry and scaling.

\item \textbf{Universality and edge behaviour.} Under broad conditions, local
eigenvalue spacings and extreme eigenvalues follow the same limiting laws
across many entry distributions.  Heavy tails and strong dependence can
invalidate this conclusion, so the hypotheses must be checked.

\end{enumerate}
\noindent\emph{Source for these results:} \cite{random_matrix_theory_ref1}.

\paragraph{Local statistics: the GUE nearest-neighbour spacing distribution}
Beyond the global eigenvalue density, random matrix theory characterizes
how individual eigenvalues are positioned relative to their neighbours. For
the Gaussian Unitary Ensemble (GUE)---the ensemble of Hermitian matrices
with independent complex Gaussian entries---the nearest-neighbour spacing
distribution $p(s)$ gives the probability density that the gap between two
adjacent eigenvalues equals $s$ after local rescaling. Unlike independent
random variables, whose gaps follow an exponential distribution, GUE
eigenvalues repel one another: $p(s)\sim(\pi s/2)e^{-\pi s^2/4}$ as
$s\to0$, so small gaps are suppressed. This level repulsion is a universal
feature of chaotic quantum systems and is experimentally observed in nuclear
spectra, mesoscopic conductance fluctuations, and zeros of the Riemann zeta
function \cite{random_matrix_theory_ref1}.

\paragraph{The Tracy--Widom distribution}
The largest eigenvalue of a large GUE matrix, after centring and rescaling,
converges to the Tracy--Widom distribution $F_2$. This distribution governs
the fluctuations of the spectral edge and appears in a remarkable range of
problems: the longest increasing subsequence of a random permutation, the
height of growth models in the Kardar--Parisi--Zhang universality class, and
the free energy of certain directed polymer models. The Tracy--Widom law
replaces the Gaussian distribution that governs sums of independent random
variables, reflecting the strong correlations induced by eigenvalue
repulsion \cite{random_matrix_theory_ref1,random_matrix_theory_ref2}.

\section*{5. Applications}
Random matrix theory studies the collective behaviour of eigenvalues and singular values when a matrix contains uncertainty. It provides reference distributions for separating stable structure from noise in large systems.
\begin{enumerate}[leftmargin=*,itemsep=0.7em]
\item \textbf{Wireless communication.}  In a multiple-input, multiple-output
(MIMO) system, a transmitter sends a vector of signals through many antenna
paths.  The receiver observes $y=Hx+z$, where $H$ is the channel matrix and
$z$ is noise.  Engineers want to know how many independent streams can be
sent reliably.  The singular values of $H$ are the strengths of those
streams, and the channel capacity contains terms of the form
$\log(1+\rho s_i^2)$, where $s_i$ are singular values and $\rho$ is the
signal-to-noise ratio.  Measuring every future channel matrix is impossible.
A random matrix model predicts the typical distribution of the $s_i$, so an
engineer can estimate capacity and the probability of a communication
outage before installing thousands of antennas.  The useful output is not
the entries of one invented matrix; it is a reliable estimate of the
distribution of possible channel qualities \cite{random_matrix_theory_ref1}.

For large antenna arrays with $n_t$ transmit and $n_r$ receive antennas,
when $n_t$ and $n_r$ grow proportionally with ratio $\gamma=n_t/n_r$, the
Marchenko--Pastur law describes the bulk of the singular-value spectrum of
$H$. The ergodic capacity per antenna approaches
\[
C=\int\log(1+\rho\,\lambda)\,dF(\lambda),
\]
where $F$ is the limiting singular-value distribution. The Tracy--Widom
distribution controls the fluctuations of the largest singular value, which
determines the dominant spatial channel and sets the peak data rate for a
single stream. These random-matrix predictions allow engineers to dimension
antenna arrays and power budgets without exhaustive field measurements,
capturing both typical performance and the probability of an unusually weak
or strong channel \cite{random_matrix_theory_ref1}.

\item \textbf{Separating signal from noise in data.}  Consider 10,000 patients
described by 1,000 laboratory measurements.  Principal component analysis
will always produce 1,000 eigenvalues, even if the measurements are
independent noise.  Random matrix theory gives a noise-only interval for
those eigenvalues.  Components inside the interval are consistent with
sampling noise; a component well outside it is a candidate for a shared
biological factor.  This is why the theory is useful: it supplies a
reference distribution against which a data scientist can test an apparent
pattern, instead of keeping the largest component merely because it is
largest.  The test must be recalibrated when measurements are strongly
dependent or have unequal variances \cite{random_matrix_theory_ref2}.

\item \textbf{Financial covariance matrices.}  A portfolio manager may estimate
the correlation of 500 assets from only 250 days of prices.  The sample
covariance matrix then has more variables than observations, and many
eigenvalues describe estimation noise rather than stable economic factors.
If the noisy covariance is inverted directly, the resulting portfolio may
place enormous weights on directions that will disappear next month.
Random matrix methods identify the noise-like bulk and replace it with a
smoothed bulk while retaining a few defensible large factors.  The resulting
matrix is better conditioned, so optimisation is less sensitive to
accidental features of the chosen historical window.  This does not predict
market prices; it improves the reliability of a matrix used by a separate
risk model.

\item \textbf{Energy levels in complex quantum systems.}  A heavy nucleus has many
interacting energy levels.  Wigner's random-matrix model predicts that nearby
levels tend to repel one another: two levels are less likely to be extremely
close than they would be in a list of independent random numbers.  Physicists
compare the rescaled spacings of measured levels with the predicted spacing
law.  Agreement is evidence that the system has the complicated mixing
expected of a chaotic quantum system; disagreement can reveal hidden
symmetry or a more regular mechanism.  The model describes statistical
regularities, not the exact energy of a particular level \cite{random_matrix_theory_ref1}.
\end{enumerate}

\theoryconnections{random_matrix_theory}{%
\item Probability supplies a distribution for the entries or eigenvalues, and linear algebra turns a matrix into eigenvalues, singular values, and eigenvectors.
\item Calculus and analysis then study the collective behavior of these quantities as the matrix size grows.
\item The important shift is from asking for one exact eigenvalue to asking for a distribution: for a large random matrix, the empirical eigenvalue plot often settles near a deterministic density, while the extreme eigenvalues fluctuate on a different scale \cite{random_matrix_theory_ref1,random_matrix_theory_ref2}.
}{%
\item In statistics, the singular values of a random data matrix provide a baseline for deciding whether a large observed component is signal or noise.
\item In wireless engineering, a channel matrix determines how many independent data streams can be transmitted, and random-matrix laws predict performance when antennas are numerous.
\item In physics, the same eigenvalue statistics model complicated energy levels, so the theory explains why unrelated systems can share the same fluctuation patterns \cite{random_matrix_theory_ref1,random_matrix_theory_ref2}.
}

\section*{Glossary}
\begin{description}[leftmargin=*,style=nextline,itemsep=0.3em]
\item[\textbf{Random matrix.}] A matrix whose entries or eigenvalue distribution are governed by a specified probability law, used to model systems with many unknown interactions.
\item[\textbf{Empirical spectral measure.}] The normalized counting distribution of eigenvalues $N^{-1}\sum_i\delta_{\lambda_i}$, approximating where eigenvalues of a large matrix tend to lie.
\item[\textbf{Marchenko--Pastur law.}] The limiting eigenvalue distribution of a large sample covariance matrix built from independent noise, defining the range of eigenvalues that noise alone would produce.
\item[\textbf{Wigner semicircle law.}] The theorem that the rescaled eigenvalue histogram of a large symmetric random matrix with centred entries approaches a semicircular shape.
\item[\textbf{Tracy--Widom distribution.}] The limiting distribution of the largest eigenvalue of certain random matrix ensembles, replacing the Gaussian in systems with strong eigenvalue correlations.
\item[\textbf{Level repulsion.}] The tendency of adjacent eigenvalues to avoid being close together, a universal feature of random matrix models and chaotic quantum systems.
\item[\textbf{Universality.}] The principle that local eigenvalue statistics and extreme eigenvalue behaviour follow the same limiting laws across broad classes of entry distributions.
\item[\textbf{Singular value.}] The square root of an eigenvalue of $M^*M$ for a matrix $M$; it measures the amplitude gain along each independent direction of a linear map.
\item[\textbf{Covariance matrix.}] A symmetric matrix whose $(i,j)$ entry is the correlation between variables $i$ and $j$.
\item[\textbf{Principal component analysis (PCA).}] A statistical technique rotating data into eigenvectors of the sample covariance matrix.
\item[\textbf{Gaussian Unitary Ensemble (GUE).}] The ensemble of $N\times N$ Hermitian matrices with independent complex Gaussian entries above the diagonal.
\item[\textbf{Sample covariance matrix.}] The matrix $S=(1/n)XX^{\mathsf T}$ computed from $n$ observations of $p$ variables.
\end{description}
\section*{7. Sample Questions}
\emph{Exercise reference:} \cite{random_matrix_theory_ref1}. The cited edition is the source for the exercise set; consult its Exercises section for the official statement and numbering.
\begin{enumerate}[leftmargin=*,itemsep=0.9em]
\item \textbf{Exercise 1.} Compute the eigenvalues of the $2\times 2$ real symmetric random matrix $M=\begin{pmatrix}1&2\\2&1\end{pmatrix}$. \emph{Solution.} The characteristic polynomial is $\det(M-\lambda I)=(1-\lambda)^2-4=\lambda^2-2\lambda-3=(\lambda-3)(\lambda+1)$. The eigenvalues are $\lambda_1=3$ and $\lambda_2=-1$.

\item \textbf{Exercise 2.} Let $X$ be a $3\times 2$ matrix with all entries i.i.d.\ standard normal $N(0,1)$. Compute the expected value of $\operatorname{Tr}(XX^T)$ and the expected sum of squared singular values of $X$. \emph{Solution.} $XX^T$ is $3\times 3$. $\operatorname{Tr}(XX^T)=\sum_{i,j}X_{ij}^2$. Each $X_{ij}^2$ has expectation $1$, and there are $3\times 2=6$ entries. So $E[\operatorname{Tr}(XX^T)]=6$. The sum of squared singular values of $X$ equals $\operatorname{Tr}(X^TX)=\operatorname{Tr}(XX^T)=6$, so the answer is $6$.

\item \textbf{Exercise 3.} For a $2\times 2$ Wigner matrix $W=\begin{pmatrix}a&b\\b&c\end{pmatrix}$ with $a,c\sim N(0,1)$ and $b\sim N(0,1/2)$ (entries independent), compute the mean and variance of $\operatorname{Tr}(W)=a+c$. \emph{Solution.} $E[\operatorname{Tr}(W)]=E[a]+E[c]=0+0=0$. $\operatorname{Var}(\operatorname{Tr}(W))=\operatorname{Var}(a)+\operatorname{Var}(c)=1+1=2$.

\item \textbf{Exercise 4.} A $4\times 4$ sample covariance matrix $S=\frac{1}{n}XX^T$ is formed from $n=10$ observations of $p=4$ variables. According to the Marchenko--Pastur law with ratio $\gamma=p/n=0.4$, what are the upper and lower edges of the noise bulk? \emph{Solution.} For the Marchenko--Pastur distribution with ratio $\gamma$, the support is $[(1-\sqrt{\gamma})^2,(1+\sqrt{\gamma})^2]$. With $\gamma=0.4$: lower edge $=(1-\sqrt{0.4})^2=(1-0.6325)^2\approx 0.1353$. Upper edge $=(1+\sqrt{0.4})^2=(1+0.6325)^2\approx 2.6649$. Eigenvalues of $S$ beyond approximately $2.66$ are candidates for signal.

\item \textbf{Exercise 5.} Let $A=\begin{pmatrix}1&0\\0&0\end{pmatrix}$. Compute the singular values of $A$ and verify that $\sum s_i^2=\operatorname{Tr}(AA^T)$. \emph{Solution.} $AA^T=A^2=\begin{pmatrix}1&0\\0&0\end{pmatrix}$. The eigenvalues of $AA^T$ are $1$ and $0$, so the singular values are $s_1=1$ and $s_2=0$. $\sum s_i^2=1^2+0^2=1$. $\operatorname{Tr}(AA^T)=1+0=1$. The identity holds.

\end{enumerate}
\section*{8. References}
\begin{thebibliography}{9}
\bibitem{random_matrix_theory_ref1} M. L. Mehta, \emph{Random Matrices}, 3rd ed., Academic Press, 2004.
\bibitem{random_matrix_theory_ref2} T. Tao, \emph{Topics in Random Matrix Theory}, American Mathematical Society, 2012.
\bibitem{random_matrix_theory_ref3} C. Tracy and H. Widom, ``Level-spacing distributions and the Airy kernel,'' \emph{Communications in Mathematical Physics}, vol.\ 159, pp.\ 151--174, 1994.
\end{thebibliography}
